%!TEX program = xelatex
\documentclass[UTF8,10pt,a4paper,fontset=none]{ctexart}

\usepackage[a4paper,top=1.1cm,bottom=1.1cm,left=1.25cm,right=1.25cm]{geometry}
\usepackage{amsmath,amssymb}
\usepackage{booktabs}
\usepackage{array}
\usepackage[most]{tcolorbox}
\usepackage{xcolor}
\usepackage{tikz}
\usetikzlibrary{arrows.meta,positioning}
\usepackage{enumitem}
\usepackage[hidelinks]{hyperref}

\setCJKmainfont{Heiti SC}
\setCJKsansfont{Heiti SC}
\setCJKmonofont{Hiragino Sans GB}

\definecolor{Ink}{HTML}{172033}
\definecolor{Blue}{HTML}{1F4E79}
\definecolor{Green}{HTML}{2A9D8F}
\definecolor{Orange}{HTML}{C86B1F}
\definecolor{SoftBlue}{HTML}{EAF3F8}
\definecolor{SoftGreen}{HTML}{EEF6EA}
\definecolor{SoftOrange}{HTML}{F8EFE4}

\setlength{\parindent}{1.5em}
\setlength{\parskip}{0.08em}
\linespread{1.02}
\setlist[itemize]{leftmargin=1.25em,itemsep=0.04em,topsep=0.05em}
\setlist[enumerate]{leftmargin=1.35em,itemsep=0.04em,topsep=0.05em}
\ctexset{section={format=\large\bfseries,beforeskip=0.45em,afterskip=0.18em}}

\newtcolorbox{note}[2][]{
  enhanced,colback=#2,colframe=Ink!65,boxrule=0.45pt,arc=0.8mm,
  left=1mm,right=1mm,top=0.7mm,bottom=0.7mm,title=#1,fonttitle=\bfseries\small
}
\tikzset{box/.style={draw,rounded corners=1pt,align=center,font=\scriptsize,minimum height=5.8mm,text width=2.05cm},arr/.style={-{Stealth[length=1.8mm]},line width=0.45pt}}

\title{\vspace{-0.9cm}\textbf{第2讲：词表示}\\[-1mm]\normalsize 单隐藏层 MLP 版朴素 CBOW 与静态词向量反思}
\author{高展沛}
\date{\vspace{-0.45cm}2026年6月}

\begin{document}
\maketitle
\vspace{-0.65cm}

\section{课堂任务：朴素 CBOW}
给定语料 $w_1,\ldots,w_T$，窗口半径 $c$，CBOW 用中心词两侧上下文
$C_t=\{w_{t-c},\ldots,w_{t-1},w_{t+1},\ldots,w_{t+c}\}$ 预测 $w_t$：
\[
\max_\theta \sum_t \log p_\theta(w_t\mid C_t).
\]
单隐藏层 MLP 结构如下：
\[
x_t=\frac{1}{|C_t|}\sum_{w\in C_t}E[w],\qquad
h_t=\tanh(W_hx_t+b_h),
\]
\[
p(w_t\mid C_t)=\mathrm{softmax}(W_oh_t+b_o),\qquad
\mathcal{L}_t=-\log p(w_t\mid C_t).
\]
其中 $E$ 是词向量矩阵，也是训练结束后要保存的表示。

\begin{center}
\resizebox{0.96\linewidth}{!}{%
\begin{tikzpicture}[node distance=5mm]
\node[box,fill=SoftBlue,draw=Blue] (a) {上下文词\\one-hot};
\node[box,fill=SoftBlue,draw=Blue,right=of a] (b) {共享\\Embedding};
\node[box,fill=SoftGreen,draw=Green,right=of b] (c) {平均池化\\$x_t$};
\node[box,fill=SoftGreen,draw=Green,right=of c] (d) {隐藏层\\$\tanh$};
\node[box,fill=SoftOrange,draw=Orange,right=of d] (e) {Softmax\\中心词};
\draw[arr] (a) -- (b);\draw[arr] (b) -- (c);\draw[arr] (c) -- (d);\draw[arr] (d) -- (e);
\end{tikzpicture}}
\end{center}

\begin{note}[手写伪代码]{SoftBlue}
\begin{enumerate}
  \item 建立词表，把语料转成 id 序列。
  \item 对每个位置 $t$，收集窗口内非中心词 id，组成训练样本 $(C_t,w_t)$。
  \item 查表得到 $E[C_t]$，对上下文向量求平均得到 $x_t$。
  \item 计算 $h_t=\tanh(W_hx_t+b_h)$ 与 $p=\mathrm{softmax}(W_oh_t+b_o)$。
  \item 用交叉熵对目标词 $w_t$ 求损失，反向传播更新 $E,W_h,b_h,W_o,b_o$。
  \item 每轮输出平均 loss；训练结束保存 $E$ 作为词表示。
\end{enumerate}
\end{note}

Skip-Gram 只把预测方向反过来：输入中心词 $w_t$，逐个预测窗口上下文 $w_j$。CBOW 训练快、对高频模式稳定；Skip-Gram 样本更多，通常对低频词更友好。

\section{探索：问题与解决}
\begin{tabular}{p{0.28\linewidth}p{0.36\linewidth}p{0.28\linewidth}}
\toprule
问题 & 原因 & 解法 \\
\midrule
完整 softmax 昂贵 & 每步归一化全词表 $V$ & negative sampling / hierarchical softmax \\
高频词支配 & 共现次数不等于信息量 & subsampling / tf-idf 式重加权 \\
OOV 与低频词差 & 词被当作不可拆原子 & fastText/subword/字符 n-gram \\
一词多义 & 一个词只有一个静态向量 & ELMo/BERT 上下文化表示 \\
长距离依赖弱 & 固定窗口只看局部 & RNN/Transformer 全句建模 \\
词序信息弱 & CBOW 平均池化丢顺序 & 位置编码、序列模型、注意力 \\
\bottomrule
\end{tabular}

\begin{note}[核心判断]{SoftGreen}
CBOW/Skip-Gram 的历史价值是把离散符号放入连续语义空间；它们的根本局限是把语义压成“静态词点”。突破方向不是简单堆深 MLP，而是让表示沿三条线进化：从局部窗口到全局统计，从词级原子到子词组合，从静态向量到上下文化动态表示。
\end{note}

\section{创新性总结}
本作业把词表示学习理解成“压缩与恢复”的过程：CBOW 把窗口上下文压缩为一个向量，再用它恢复中心词。若恢复失败，错误通常不是某个词孤立造成的，而是采样、词表、语义歧义和上下文范围共同造成的。因此，一个更现代的路线应是：先用负采样和高频词控制提高可训练性，再用 subword 解决泛化，最后用 Transformer/RNN 让同一个词在不同语境中拥有不同表示。这也自然引出第5讲的序列标注：词向量只是起点，真正的实体判断必须回到上下文中完成。

\end{document}
